\hypertarget{the-policy-language}{%
\chapter{20. The policy language}\label{the-policy-language}}

\emph{Generated from \texttt{schema/policy.toml.schema}. Current as of
0.5.0.}

\hypertarget{audit}{%
\section{\texorpdfstring{\texttt{{[}audit{]}}}{{[}audit{]}}}\label{audit}}

\texttt{{[}audit{]}}: sink selection, per-class levels, and per-sink
tuning

\begin{kennelcode}
[audit]
# optional
sinks    = ["..."]
\end{kennelcode}

Contains: \texttt{{[}audit.dbus{]}}, \texttt{{[}audit.exec{]}},
\texttt{{[}audit.file{]}}, \texttt{{[}audit.filesystem{]}},
\texttt{{[}audit.journald{]}}, \texttt{{[}audit.network{]}},
\texttt{{[}audit.stdout{]}}, \texttt{{[}audit.syslog{]}},
\texttt{{[}audit.unix{]}}.

\begin{description}
\tightlist
\item[\texttt{sinks}]
optional. Active sinks (\texttt{file}, \texttt{journald},
\texttt{syslog}, \texttt{stdout}). Default \texttt{{[}"file"{]}}.
\end{description}

\hypertarget{audit.dbus}{%
\subsection{\texorpdfstring{\texttt{{[}audit.dbus{]}}}{{[}audit.dbus{]}}}\label{audit.dbus}}

\texttt{{[}audit.dbus{]}} level.

\begin{kennelcode}
[audit.dbus]
# optional
level    = "off" | "denies-only" | "summary" | "full"
\end{kennelcode}

\begin{description}
\tightlist
\item[\texttt{level}]
optional. One of \texttt{off}, \texttt{denies-only}, \texttt{summary},
\texttt{full}. One of \texttt{off}, \texttt{denies-only},
\texttt{summary}, \texttt{full}.
\end{description}

\hypertarget{audit.file}{%
\subsection{\texorpdfstring{\texttt{{[}audit.file{]}}}{{[}audit.file{]}}}\label{audit.file}}

\texttt{{[}audit.file{]}} tuning.

\begin{kennelcode}
[audit.file]
# optional
compress_after_seconds = 0
dir      = "..."
retain_count = 0
rotate_at_bytes = "..."
\end{kennelcode}

\begin{description}
\tightlist
\item[\texttt{compress\_after\_seconds}]
optional. Gzip a rotated file this many seconds after rotation.
\item[\texttt{dir}]
optional. Override the per-kennel directory (placeholders allowed).
\item[\texttt{retain\_count}]
optional. Keep at most this many rotated files per class.
\item[\texttt{rotate\_at\_bytes}]
optional. Rotate at this size (e.g.~\texttt{"64M"}, \texttt{"1G"}; bare
= bytes).
\end{description}

\hypertarget{audit.journald}{%
\subsection{\texorpdfstring{\texttt{{[}audit.journald{]}}}{{[}audit.journald{]}}}\label{audit.journald}}

\texttt{{[}audit.journald{]}} --- no fields; present to allow the empty
table.

\hypertarget{audit.syslog}{%
\subsection{\texorpdfstring{\texttt{{[}audit.syslog{]}}}{{[}audit.syslog{]}}}\label{audit.syslog}}

\texttt{{[}audit.syslog{]}} tuning.

\begin{kennelcode}
[audit.syslog]
# optional
facility = "..."
\end{kennelcode}

\begin{description}
\tightlist
\item[\texttt{facility}]
optional. Syslog facility (\texttt{user}, \texttt{daemon},
\texttt{auth}, \ldots). Default \texttt{user}.
\end{description}

\hypertarget{cap}{%
\section{\texorpdfstring{\texttt{{[}cap{]}}}{{[}cap{]}}}\label{cap}}

\texttt{{[}cap{]}}: capabilities and \texttt{no\_new\_privs}.

\begin{kennelcode}
[cap]
# optional
bounding_set = ["..."]
no_new_privs = true
\end{kennelcode}

\begin{description}
\tightlist
\item[\texttt{bounding\_set}]
optional. The capability bounding set to retain (empty drops them all).
\item[\texttt{no\_new\_privs}]
optional. \texttt{PR\_SET\_NO\_NEW\_PRIVS}. A framework invariant once
resolved (must be true).
\end{description}

Example, from \texttt{toml/templates/base-bwrap}:

\begin{kennelcode}
[cap]
no_new_privs = true
bounding_set = []
\end{kennelcode}

\hypertarget{consumes}{%
\section{\texorpdfstring{\texttt{{[}{[}consumes{]}{]}}}{{[}{[}consumes{]}{]}}}\label{consumes}}

One \texttt{{[}{[}consumes{]}{]}} entry, a capability this kennel
reaches over the mesh.

\begin{kennelcode}
[[consumes]]
# required
name     = "..."
shape    = "af-unix" | "dbus-name" | "binder-connector"
# optional
at       = "..."
env      = ["..."]
key      = "..."
required = true
# decoration
reason   = "..."
\end{kennelcode}

\begin{description}
\tightlist
\item[\texttt{at}]
optional. Where the brokered connector is delivered, in this kennel's
own view.
\item[\texttt{env}]
optional. Environment variable(s) synthesised into this kennel to name
the connector.
\item[\texttt{key}]
optional. An optional private match token; must match the provider's.
\item[\texttt{name}]
required. The capability's public identifier, resolved against the
catalogue at runtime.
\item[\texttt{reason}]
decoration. Why this capability is consumed (required).
\item[\texttt{required}]
optional. Whether the capability's absence fails kennel construction.
Hard dependency by
\item[\texttt{shape}]
required. The transport it expects; the broker refuses a mismatched
shape. One of \texttt{af-unix}, \texttt{dbus-name},
\texttt{binder-connector}.
\end{description}

Example, from \texttt{toml/templates/gui-session}:

\begin{kennelcode}
[[consumes]]
name = "org.projectkennel.wayland"
shape = "af-unix"
at = "/tmp/wayland-0"
required = true
reason = "the confined desktop session reaches its display server — a broker-spawned nested compositor — over the mesh"
\end{kennelcode}

\hypertarget{dbus}{%
\section{\texorpdfstring{\texttt{{[}dbus{]}}}{{[}dbus{]}}}\label{dbus}}

\texttt{{[}dbus{]}}: D-Bus mediation.

Contains: \texttt{{[}dbus.audit{]}}, \texttt{{[}dbus.session{]}},
\texttt{{[}dbus.system{]}}.

\hypertarget{dbus.audit}{%
\subsection{\texorpdfstring{\texttt{{[}dbus.audit{]}}}{{[}dbus.audit{]}}}\label{dbus.audit}}

\texttt{{[}dbus.audit{]}}: per-kennel D-Bus call audit verbosity.

\begin{kennelcode}
[dbus.audit]
# optional
level    = "off" | "summary" | "full"
\end{kennelcode}

\begin{description}
\tightlist
\item[\texttt{level}]
optional. Verbosity (\texttt{"off"}, \texttt{"summary"},
\texttt{"full"}). One of \texttt{off}, \texttt{summary}, \texttt{full}.
\end{description}

\hypertarget{dbus.session}{%
\subsection{\texorpdfstring{\texttt{{[}dbus.session{]}}}{{[}dbus.session{]}}}\label{dbus.session}}

\texttt{{[}dbus.session{]}}: the user session bus (the common case:
notifications, portals).

\begin{kennelcode}
[dbus.session]
# optional
enabled  = true
\end{kennelcode}

Contains: \texttt{{[}dbus.session.allow{]}},
\texttt{{[}dbus.session.deny{]}}.

\begin{description}
\tightlist
\item[\texttt{enabled}]
optional. Whether this bus is reachable at all. Absent/\texttt{false} ⇒
no connection to this bus.
\end{description}

\hypertarget{dbus.session.allow}{%
\subsubsection{\texorpdfstring{\texttt{{[}dbus.session.allow{]}}}{{[}dbus.session.allow{]}}}\label{dbus.session.allow}}

\texttt{{[}dbus.\textless{}bus\textgreater{}.allow{]}}: what the kennel
may reach (an allowlist; default-deny).

\begin{kennelcode}
[dbus.session.allow]
# optional
broadcast = ["..."]
call     = ["..."]
own      = ["..."]
talk     = ["..."]
\end{kennelcode}

\begin{description}
\tightlist
\item[\texttt{broadcast}]
optional. Signals the kennel may receive (a subset of senders it may
\texttt{talk} to).
\item[\texttt{call}]
optional. Finer than \texttt{talk}: specific
\texttt{destination=interface.member} calls.
\item[\texttt{own}]
optional. Names the kennel may own (be addressable as). Almost always
empty.
\item[\texttt{talk}]
optional. Destinations the kennel may call methods on and receive
replies/signals from
\end{description}

\hypertarget{env}{%
\section{\texorpdfstring{\texttt{{[}env{]}}}{{[}env{]}}}\label{env}}

\texttt{{[}env{]}}: environment curation.

\begin{kennelcode}
[env]
# optional
deny     = ["..."]
pass     = ["..."]
set      = "..."
\end{kennelcode}

\begin{description}
\tightlist
\item[\texttt{deny}]
optional. Variables denied even if passed (globs allowed).
\item[\texttt{pass}]
optional. Variables passed through from the caller's environment (globs
allowed).
\item[\texttt{set}]
optional. Variables forced to a specific value. Declared last: as a TOML
table it must
\end{description}

Example, from \texttt{toml/templates/base-bwrap}:

\begin{kennelcode}
[env]
set = { TMPDIR = "/tmp", NO_PROXY = "" }
\end{kennelcode}

\hypertarget{exec}{%
\section{\texorpdfstring{\texttt{{[}exec{]}}}{{[}exec{]}}}\label{exec}}

\texttt{{[}exec{]}}: what may be \texttt{execve}'d.

\begin{kennelcode}
[exec]
# optional
allow    = ["..."]
deny     = ["..."]
deny_setcap = true
deny_setgid = true
deny_setuid = true
deny_writable = true
path     = ["..."]
shell    = "..."
\end{kennelcode}

\begin{description}
\tightlist
\item[\texttt{allow}]
optional. Allowlisted binary paths (the execve allowlist). Execution is
deny-by-default:
\item[\texttt{deny}]
optional. Denylisted absolute paths or globs.
\item[\texttt{deny\_setcap}]
optional. Refuse file-capability binaries (framework invariant).
\item[\texttt{deny\_setgid}]
optional. Refuse setgid binaries (framework invariant).
\item[\texttt{deny\_setuid}]
optional. Refuse setuid binaries at execve (framework invariant once
resolved).
\item[\texttt{deny\_writable}]
optional. Refuse execution of files in writable paths (framework
invariant).
\item[\texttt{path}]
optional. \texttt{PATH} search roots the resolver records for the
workload's environment.
\item[\texttt{shell}]
optional. The kennel's login shell: the synthetic-\texttt{passwd}
\texttt{pw\_shell} and
\end{description}

Example, from \texttt{toml/templates/argv-tool}:

\begin{kennelcode}
[exec]
allow = ["/bin/echo", "/usr/bin/echo", "/bin/cat", "/usr/bin/cat", "/bin/true", "/usr/bin/true", "/bin/sh"]
\end{kennelcode}

\hypertarget{fs}{%
\section{\texorpdfstring{\texttt{{[}fs{]}}}{{[}fs{]}}}\label{fs}}

\texttt{{[}fs{]}} and its sub-tables.

\begin{kennelcode}
[fs]
# optional
deny     = ["..."]
exclusive = ["..."]
read     = ["..."]
write    = ["..."]
\end{kennelcode}

Contains: \texttt{{[}fs.dev{]}}, \texttt{{[}fs.home{]}},
\texttt{{[}fs.proc{]}}, \texttt{{[}fs.tmp{]}}.

\begin{description}
\tightlist
\item[\texttt{deny}]
optional. Categorical denies (belt-and-braces over the constructed
view). Replace or increment
\item[\texttt{exclusive}]
optional. Writable paths bound \textbf{exclusively} (T2.8): while the
kennel runs, \texttt{kenneld}
\item[\texttt{read}]
optional. Paths granted read (and directory traversal / execute).
Replace (\texttt{read\ =\ {[}"…"{]}}) or
\item[\texttt{write}]
optional. Paths granted write. Replace or increment at the same key.
\end{description}

Example, from \texttt{toml/templates/argv-tool}:

\begin{kennelcode}
[fs]
read = ["/usr/**", "/bin/**", "/lib/**", "/lib64/**"]
\end{kennelcode}

\hypertarget{fs.dev}{%
\subsection{\texorpdfstring{\texttt{{[}fs.dev{]}}}{{[}fs.dev{]}}}\label{fs.dev}}

\texttt{{[}fs.dev{]}}: the minimal \texttt{/dev}.

\begin{kennelcode}
[fs.dev]
# optional
allow    = ["..."]
\end{kennelcode}

Contains: \texttt{{[}{[}fs.dev.passthrough{]}{]}}.

\begin{description}
\tightlist
\item[\texttt{allow}]
optional. The trivial pseudo-device baseline bound into the kennel's
\texttt{/dev} (\texttt{/dev/null},
\end{description}

\textbf{\texttt{{[}{[}fs.dev.passthrough{]}{]}}} entries (or
\texttt{\{\ add,\ remove\ \}} increment), optional:

\begin{description}
\tightlist
\item[\texttt{group}]
optional. The owning group that gates access (e.g.~\texttt{dialout},
\texttt{modem}, \texttt{dip}). Access is
\item[\texttt{path}]
optional. The device node, an absolute path under \texttt{/dev}
(e.g.~\texttt{/dev/ttyUSB0},
\item[\texttt{reason}]
decoration. Why this device is exposed (required).
\item[\texttt{threats}]
decoration. Threat tags, required to carry an \texttt{exposed} tag
(passthrough widens the
\end{description}

Example, from \texttt{toml/templates/base-bwrap}:

\begin{kennelcode}
[fs.dev]
allow = [
    "/dev/null", "/dev/zero", "/dev/random", "/dev/urandom",
    "/dev/tty", "/dev/pts/**",
]
\end{kennelcode}

\hypertarget{fs.home}{%
\subsection{\texorpdfstring{\texttt{{[}fs.home{]}}}{{[}fs.home{]}}}\label{fs.home}}

\texttt{{[}fs.home{]}}: the constructed \texttt{\$HOME} view.

\begin{kennelcode}
[fs.home]
# optional
persist  = ["..."]
readonly = true
shadow   = true
\end{kennelcode}

\begin{description}
\tightlist
\item[\texttt{persist}]
optional. Home-relative paths that \textbf{persist} across runs. By
default the
\item[\texttt{readonly}]
optional. Make the constructed \texttt{\$HOME} \textbf{read-only}
(default: writable). The home root
\item[\texttt{shadow}]
optional. Whether \texttt{\$HOME} is shadowed by a constructed view
(must be true once resolved).
\end{description}

Example, from \texttt{toml/templates/base-bwrap}:

\begin{kennelcode}
[fs.home]
shadow = true
\end{kennelcode}

\hypertarget{fs.proc}{%
\subsection{\texorpdfstring{\texttt{{[}fs.proc{]}}}{{[}fs.proc{]}}}\label{fs.proc}}

\texttt{{[}fs.proc{]}}: procfs hidepid.

\begin{kennelcode}
[fs.proc]
# optional
hidepid  = true
\end{kennelcode}

\begin{description}
\tightlist
\item[\texttt{hidepid}]
optional. Mount \texttt{/proc} with \texttt{hidepid=2}.
\end{description}

Example, from \texttt{toml/templates/base-bwrap}:

\begin{kennelcode}
[fs.proc]
hidepid = true
\end{kennelcode}

\hypertarget{fs.tmp}{%
\subsection{\texorpdfstring{\texttt{{[}fs.tmp{]}}}{{[}fs.tmp{]}}}\label{fs.tmp}}

\texttt{{[}fs.tmp{]}}: the private \texttt{/tmp} tmpfs.

\begin{kennelcode}
[fs.tmp]
# optional
size     = "..."
writable = true
\end{kennelcode}

\begin{description}
\tightlist
\item[\texttt{size}]
optional. Size cap in human form (\texttt{"512M"}, \texttt{"1G"}).
\item[\texttt{writable}]
optional. Whether the workload may \textbf{write} to its \texttt{/tmp}
tmpfs (the Landlock write grant). Absent ⇒
\end{description}

Example, from \texttt{toml/templates/base-bwrap}:

\begin{kennelcode}
[fs.tmp]
writable = true
size = "512M"
\end{kennelcode}

\hypertarget{identity}{%
\section{\texorpdfstring{\texttt{{[}identity{]}}}{{[}identity{]}}}\label{identity}}

\texttt{{[}identity{]}}: the workload's identity inside the kennel.

\begin{kennelcode}
[identity]
# optional
group    = "..."
groups   = ["..."]
user     = "..."
\end{kennelcode}

\begin{description}
\tightlist
\item[\texttt{group}]
optional. The workload's masked \textbf{primary} group name (synthetic
\texttt{/etc/passwd} \texttt{pw\_gid}
\item[\texttt{groups}]
optional. Supplementary group names to retain
(e.g.~\texttt{{[}"dialout",\ "plugdev"{]}}). The user
\item[\texttt{user}]
optional. The workload's masked user name,
\texttt{\$USER}/\texttt{\$LOGNAME} and the synthetic
\end{description}

\hypertarget{lifecycle}{%
\section{\texorpdfstring{\texttt{{[}lifecycle{]}}}{{[}lifecycle{]}}}\label{lifecycle}}

\texttt{{[}lifecycle{]}}: TTL and TTL action. \texttt{ttl} is the human
form (\texttt{"8h"}); the

\begin{kennelcode}
[lifecycle]
# optional
ttl      = "..."
ttl_action = "exit" | "warn" | "renew"
\end{kennelcode}

\begin{description}
\tightlist
\item[\texttt{ttl}]
optional. Time-to-live in human form (\texttt{"8h"}, \texttt{"1h"},
\texttt{"30m"}).
\item[\texttt{ttl\_action}]
optional. What to do at TTL expiry: \texttt{"exit"} (alias
\texttt{"stop"}, the default) ends the One of \texttt{exit},
\texttt{warn}, \texttt{renew}.
\end{description}

Example, from \texttt{toml/templates/argv-tool}:

\begin{kennelcode}
[lifecycle]
ttl = "5m"
ttl_action = "exit"
\end{kennelcode}

\hypertarget{mutable}{%
\section{\texorpdfstring{\texttt{{[}{[}mutable{]}{]}}}{{[}{[}mutable{]}{]}}}\label{mutable}}

One \texttt{{[}{[}mutable{]}{]}} manifest entry: a leaf field a spawn of
this template may write.

\begin{kennelcode}
[[mutable]]
# optional
field    = "..."
freeform = true
from     = ["..."]
match    = ["..."]
max      = 0
oneof    = ["..."]
type     = "..."
under    = "..."
# decoration
reason   = "..."
\end{kennelcode}

\begin{description}
\tightlist
\item[\texttt{field}]
optional. The dotted leaf-field path this entry opens
(\texttt{net.proxy.allow}, \texttt{rootfs.writable}, \texttt{fs.write}).
\item[\texttt{freeform}]
optional. Freeform bound: no shape at all, the loud last-resort footgun.
Any value is accepted; a
\item[\texttt{from}]
optional. Pool bound: the fixed set a spawn may append values from.
\item[\texttt{match}]
optional. Pattern bound: the pre-baked net-destination shapes an open
value must match
\item[\texttt{max}]
optional. Pool bound: the maximum number of appended entries.
\item[\texttt{oneof}]
optional. Oneof bound: the enumerated member list a spawn selects from.
\item[\texttt{reason}]
decoration. The justification a \texttt{freeform} variant requires (the
loud rule).
\item[\texttt{type}]
optional. Predicate bound: the value type (currently \texttt{relpath}).
\item[\texttt{under}]
optional. Predicate bound: the root the value resolves under
(\texttt{RESOLVE\_IN\_ROOT}, traversal-free).
\end{description}

Example, from \texttt{toml/templates/argv-tool}:

\begin{kennelcode}
[[mutable]]
field = "workload.argv"
freeform = true
reason = "the caller chooses the command to run within the [exec].allow floor; the cage contains it"
\end{kennelcode}

\hypertarget{net}{%
\section{\texorpdfstring{\texttt{{[}net{]}}}{{[}net{]}}}\label{net}}

\texttt{{[}net{]}} and its sub-tables.

\begin{kennelcode}
[net]
# optional
mode     = "none" | "constrained" | "unconstrained" | "host"
proxy_listen_v4_address = "..."
proxy_listen_v6_address = "..."
# decoration
reason   = "..."
\end{kennelcode}

Contains: \texttt{{[}net.audit{]}}, \texttt{{[}net.bind{]}},
\texttt{{[}net.bpf{]}}, \texttt{{[}net.ipv6{]}},
\texttt{{[}net.proxy{]}}.

\begin{description}
\tightlist
\item[\texttt{mode}]
optional. Egress mode: \texttt{"none"} (own empty net-ns, no
interfaces), \texttt{"constrained"} (own net-ns, One of \texttt{none},
\texttt{constrained}, \texttt{unconstrained}, \texttt{host}.
\item[\texttt{proxy\_listen\_v4\_address}]
optional. IPv4 proxy listen address as \texttt{"offset:port"} within the
kennel's subnet. A
\item[\texttt{proxy\_listen\_v6\_address}]
optional. IPv6 proxy listen address as \texttt{"offset:port"}.
\item[\texttt{reason}]
decoration. Required (non-empty) only when \texttt{mode\ =\ "host"}: the
documented justification for
\end{description}

Example, from \texttt{toml/policies/interactive}:

\begin{kennelcode}
[net]
mode = "host"
reason = "an interactive human shell needs direct, unmediated egress (git/ssh/curl) on the host network stack"
\end{kennelcode}

\hypertarget{net.audit}{%
\subsection{\texorpdfstring{\texttt{{[}net.audit{]}}}{{[}net.audit{]}}}\label{net.audit}}

\texttt{{[}net.audit{]}}: per-kennel egress audit log.

\begin{kennelcode}
[net.audit]
# optional
level    = "summary" | "full"
log_path = "..."
\end{kennelcode}

\begin{description}
\tightlist
\item[\texttt{level}]
optional. Audit verbosity (\texttt{"summary"}, \texttt{"full"}). One of
\texttt{summary}, \texttt{full}.
\item[\texttt{log\_path}]
optional. Where the per-kennel egress JSONL log is written.
\end{description}

Example, from \texttt{toml/templates/base-confined}:

\begin{kennelcode}
[net.audit]
log_path = "~/.local/state/kennel/<kennel>/network.jsonl"
level = "summary"
\end{kennelcode}

\hypertarget{net.bind}{%
\subsection{\texorpdfstring{\texttt{{[}net.bind{]}}}{{[}net.bind{]}}}\label{net.bind}}

\texttt{{[}net.bind{]}}: bind-address rewriting policy (the
wildcard-rewrite knobs; the bind

\begin{kennelcode}
[net.bind]
# optional
allow_host_loopback_v4 = true
allow_host_loopback_v6 = true
allowed_ports = ["..."]
in6addr_any_policy = "rewrite" | "deny"
inaddr_any_policy = "rewrite" | "deny"
min_port = 0
\end{kennelcode}

\begin{description}
\tightlist
\item[\texttt{allow\_host\_loopback\_v4}]
optional. Whether binding the host IPv4 loopback is permitted.
\item[\texttt{allow\_host\_loopback\_v6}]
optional. Whether binding the host IPv6 loopback is permitted.
\item[\texttt{allowed\_ports}]
optional. Explicit allowlist of bindable ports. When non-empty, the
workload may
\item[\texttt{in6addr\_any\_policy}]
optional. What to do with a wildcard IPv6 bind. One of \texttt{rewrite},
\texttt{deny}.
\item[\texttt{inaddr\_any\_policy}]
optional. What to do with a wildcard IPv4 bind (\texttt{"rewrite"} /
\texttt{"deny"}). One of \texttt{rewrite}, \texttt{deny}.
\item[\texttt{min\_port}]
optional. Lowest bindable port.
\end{description}

Example, from \texttt{toml/templates/base-confined}:

\begin{kennelcode}
[net.bind]
inaddr_any_policy = "rewrite"
in6addr_any_policy = "rewrite"
allow_host_loopback_v4 = false
allow_host_loopback_v6 = false
min_port = 1024
\end{kennelcode}

\hypertarget{net.bpf}{%
\subsection{\texorpdfstring{\texttt{{[}net.bpf{]}}}{{[}net.bpf{]}}}\label{net.bpf}}

\texttt{{[}net.bpf{]}}: the kernel/syscall ACL (the cgroup
\texttt{connect4/6} + \texttt{bind4/6} BPF and the

\begin{kennelcode}
[net.bpf]
# optional
deny_families = ["..."]
families = ["..."]
\end{kennelcode}

Contains: \texttt{{[}net.bpf.bind{]}}, \texttt{{[}net.bpf.connect{]}}.

\begin{description}
\tightlist
\item[\texttt{deny\_families}]
optional. Denied socket families (\texttt{inet\_sock\_create} returns
EPERM): \texttt{AF\_NETLINK}, \texttt{AF\_PACKET}, \ldots{}
\item[\texttt{families}]
optional. Permitted socket families (defence in depth;
e.g.~\texttt{{[}"AF\_INET",\ "AF\_INET6",\ "AF\_UNIX"{]}}).
\end{description}

\hypertarget{net.bpf.bind}{%
\subsubsection{\texorpdfstring{\texttt{{[}net.bpf.bind{]}}}{{[}net.bpf.bind{]}}}\label{net.bpf.bind}}

\texttt{{[}net.bpf.bind{]}}: the inbound BIND ACL (cidr + ports,
deny-first).

Contains: \texttt{{[}{[}net.bpf.bind.allow{]}{]}},
\texttt{{[}{[}net.bpf.bind.deny{]}{]}}.

\textbf{\texttt{{[}{[}net.bpf.bind.allow{]}{]}}} entries (or
\texttt{\{\ add,\ remove\ \}} increment), optional:

\begin{description}
\tightlist
\item[\texttt{cidr}]
optional. The CIDR (\texttt{"10.0.0.0/8"}, a bare address, or
\texttt{"*"} = \texttt{0.0.0.0/0} + \texttt{::/0}).
\item[\texttt{ports}]
optional. Permitted ports (empty = any port).
\item[\texttt{protocol}]
optional. Transport protocol (\texttt{"tcp"}, \texttt{"udp"},
\texttt{"any"}). One of \texttt{any}, \texttt{tcp}, \texttt{udp}.
\item[\texttt{reason}]
decoration. Why this rule exists (required).
\item[\texttt{threats}]
decoration. Threat tags.
\end{description}

\textbf{\texttt{{[}{[}net.bpf.bind.deny{]}{]}}} entries (or
\texttt{\{\ add,\ remove\ \}} increment), optional:

\begin{description}
\tightlist
\item[\texttt{cidr}]
optional. The CIDR (\texttt{"10.0.0.0/8"}, a bare address, or
\texttt{"*"} = \texttt{0.0.0.0/0} + \texttt{::/0}).
\item[\texttt{ports}]
optional. Permitted ports (empty = any port).
\item[\texttt{protocol}]
optional. Transport protocol (\texttt{"tcp"}, \texttt{"udp"},
\texttt{"any"}). One of \texttt{any}, \texttt{tcp}, \texttt{udp}.
\item[\texttt{reason}]
decoration. Why this rule exists (required).
\item[\texttt{threats}]
decoration. Threat tags.
\end{description}

\hypertarget{net.ipv6}{%
\subsection{\texorpdfstring{\texttt{{[}net.ipv6{]}}}{{[}net.ipv6{]}}}\label{net.ipv6}}

\texttt{{[}net.ipv6{]}}: IPv6-specific options.

\begin{kennelcode}
[net.ipv6]
# optional
force_v6only = true
\end{kennelcode}

\begin{description}
\tightlist
\item[\texttt{force\_v6only}]
optional. Force \texttt{IPV6\_V6ONLY=1} so a dual-stack socket cannot
escape the v4 rewrite.
\end{description}

Example, from \texttt{toml/templates/base-confined}:

\begin{kennelcode}
[net.ipv6]
force_v6only = true
\end{kennelcode}

\hypertarget{net.proxy}{%
\subsection{\texorpdfstring{\texttt{{[}net.proxy{]}}}{{[}net.proxy{]}}}\label{net.proxy}}

\texttt{{[}net.proxy{]}}: the user-space egress policy the per-kennel
proxy enforces

Contains: \texttt{{[}net.proxy.deny{]}},
\texttt{{[}{[}net.proxy.allow{]}{]}}.

\textbf{\texttt{{[}{[}net.proxy.allow{]}{]}}} entries (or
\texttt{\{\ add,\ remove\ \}} increment), optional:

\begin{description}
\tightlist
\item[\texttt{cidr}]
optional. A CIDR destination, when the rule is by-address rather than
by-name.
\item[\texttt{name}]
optional. The destination host (or dot-prefixed suffix). Mutually
informative with \texttt{cidr}.
\item[\texttt{ports}]
optional. Permitted ports.
\item[\texttt{protocol}]
optional. Transport protocol (\texttt{"tcp"}, \texttt{"udp"},
\texttt{"any"}). One of \texttt{any}, \texttt{tcp}, \texttt{udp}.
\item[\texttt{reason}]
decoration. Why this destination is permitted (required).
\item[\texttt{threats}]
decoration. Threat tags.
\item[\texttt{tls}]
optional. \texttt{tls.required} and friends.
\end{description}

\hypertarget{net.proxy.deny}{%
\subsubsection{\texorpdfstring{\texttt{{[}net.proxy.deny{]}}}{{[}net.proxy.deny{]}}}\label{net.proxy.deny}}

\texttt{{[}net.proxy.deny{]}}: the deny table: the non-removable
\texttt{invariant} floor and the

Contains: \texttt{{[}{[}net.proxy.deny.invariant{]}{]}},
\texttt{{[}{[}net.proxy.deny.policy{]}{]}}.

\textbf{\texttt{{[}{[}net.proxy.deny.invariant{]}{]}}} entries,
optional:

\begin{description}
\tightlist
\item[\texttt{cidr}]
required. The denied CIDR (e.g.~\texttt{"169.254.169.254/32"}).
\item[\texttt{reason}]
decoration. Why the deny exists (required).
\item[\texttt{threats}]
decoration. Threat tags.
\end{description}

\textbf{\texttt{{[}{[}net.proxy.deny.policy{]}{]}}} entries (or
\texttt{\{\ add,\ remove\ \}} increment), optional:

\begin{description}
\tightlist
\item[\texttt{cidr}]
required. The denied CIDR (e.g.~\texttt{"169.254.169.254/32"}).
\item[\texttt{reason}]
decoration. Why the deny exists (required).
\item[\texttt{threats}]
decoration. Threat tags.
\end{description}

\hypertarget{provides}{%
\section{\texorpdfstring{\texttt{{[}{[}provides{]}{]}}}{{[}{[}provides{]}{]}}}\label{provides}}

One \texttt{{[}{[}provides{]}{]}} entry, a capability this kennel offers
over the mesh.

\begin{kennelcode}
[[provides]]
# required
name     = "..."
shape    = "af-unix" | "dbus-name" | "binder-connector"
# optional
endpoint = "..."
key      = "..."
# decoration
reason   = "..."
\end{kennelcode}

\begin{description}
\tightlist
\item[\texttt{endpoint}]
optional. Where the capability is exposed, in the provider's own view.
Optional: an omitted
\item[\texttt{key}]
optional. An optional private match token, never advertised in the
catalogue.
\item[\texttt{name}]
required. The capability's public identifier, what the catalogue
advertises. A reserved
\item[\texttt{reason}]
decoration. Why this capability is offered (required).
\item[\texttt{shape}]
required. The typed transport. One of \texttt{af-unix},
\texttt{dbus-name}, \texttt{binder-connector}.
\end{description}

Example, from \texttt{toml/templates/dbus-broker}:

\begin{kennelcode}
[[provides]]
name = "org.projectkennel.dbus-broker"
shape = "binder-connector"
endpoint = "/dev/binderfs-mesh/binder"
reason = "kenneld pushes per-consumer ACCEPT_SESSION (bus + method filter) to the broker's control node on the mesh"
\end{kennelcode}

\hypertarget{rootfs}{%
\section{\texorpdfstring{\texttt{{[}rootfs{]}}}{{[}rootfs{]}}}\label{rootfs}}

\texttt{{[}rootfs{]}}: an OCI image unpacked as the kennel's root
filesystem.

\begin{kennelcode}
[rootfs]
# optional
image    = "..."
path     = "..."
persistence = "discard" | "persist"
readonly = ["..."]
writable = ["..."]
# decoration
reason   = "..."
\end{kennelcode}

\begin{description}
\tightlist
\item[\texttt{image}]
optional. The \texttt{image@sha256:…} the build pulled from; the runner
refuses unless it equals the
\item[\texttt{path}]
optional. The unpacked image rootfs (the store entry's
\texttt{rootfs/}).
\item[\texttt{persistence}]
optional. Rootfs persistence: \texttt{"discard"} (default) \textbar{}
\texttt{"persist"}. \texttt{"persist"} is a One of \texttt{discard},
\texttt{persist}.
\item[\texttt{readonly}]
optional. Closure-lock: rootfs paths Landlock denies writes to, the
executable-closure
\item[\texttt{reason}]
decoration. Why this substrate is trusted (required; the substrate-trust
waiver is loud).
\item[\texttt{writable}]
optional. Closure-lock holes: rootfs paths kept writable, carved back
out of \texttt{readonly}
\end{description}

\hypertarget{seccomp}{%
\section{\texorpdfstring{\texttt{{[}seccomp{]}}}{{[}seccomp{]}}}\label{seccomp}}

\texttt{{[}seccomp{]}}: the seccomp filter (source carries a deny list;
the resolver

\begin{kennelcode}
[seccomp]
# optional
allow    = ["..."]
deny     = ["..."]
profile  = "..."
\end{kennelcode}

\begin{description}
\tightlist
\item[\texttt{allow}]
optional. Syscalls explicitly allowed.
\item[\texttt{deny}]
optional. Syscalls denied on top of the profile.
\item[\texttt{profile}]
optional. The baseline profile name (\texttt{"default"}).
\end{description}

Example, from \texttt{toml/templates/base-bwrap}:

\begin{kennelcode}
[seccomp]
profile = "default"
deny = [
    "userfaultfd", "perf_event_open", "bpf",
    "process_vm_readv", "process_vm_writev",
    "kexec_load", "kexec_file_load",
    "mount", "umount", "umount2", "pivot_root",
    "swapon", "swapoff", "reboot",
    "init_module", "finit_module", "delete_module",
    "personality",
]
\end{kennelcode}

\hypertarget{service}{%
\section{\texorpdfstring{\texttt{{[}service{]}}}{{[}service{]}}}\label{service}}

\texttt{{[}service{]}}: the supervision discipline for a service kennel.

\begin{kennelcode}
[service]
# optional
backoff  = "..."
max_attempts = 0
restart  = "always" | "on-failure" | "never"
\end{kennelcode}

\begin{description}
\tightlist
\item[\texttt{backoff}]
optional. Initial delay before a restart in human form
(\texttt{"500ms"}, \texttt{"2s"}, default \texttt{"500ms"}); doubles
\item[\texttt{max\_attempts}]
optional. Restarts within the crash-loop window before
declared-but-failed (default \texttt{5}; must be ≥ 1).
\item[\texttt{restart}]
optional. Restart discipline: \texttt{always} / \texttt{on-failure}
(default) / \texttt{never}. One of \texttt{always}, \texttt{on-failure},
\texttt{never}.
\end{description}

Example, from \texttt{toml/templates/dbus-broker}:

\begin{kennelcode}
[service]
restart = "on-failure"
\end{kennelcode}

\hypertarget{signature}{%
\section{\texorpdfstring{\texttt{{[}signature{]}}}{{[}signature{]}}}\label{signature}}

The \texttt{{[}signature{]}} envelope carried by a signed artefact.

\begin{kennelcode}
[signature]
# required
algorithm = "..."
key_id   = "..."
signature = "..."
# optional
signed_fields = ["..."]
\end{kennelcode}

\begin{description}
\tightlist
\item[\texttt{algorithm}]
required. Signature algorithm. Must be \texttt{"sshsig"}
({[}\texttt{SSHSIG\_ALGORITHM}{]}).
\item[\texttt{key\_id}]
required. Identifies the signing key in the trust store. The SSHSIG also
embeds the
\item[\texttt{signature}]
required. The armored SSHSIG
(\texttt{-\/-\/-\/-\/-BEGIN\ SSH\ SIGNATURE-\/-\/-\/-\/-} \ldots) over
the canonical
\item[\texttt{signed\_fields}]
optional. The top-level fields the signature covers (every field except
\end{description}

Example, from \texttt{toml/policies/interactive}:

\begin{kennelcode}
[signature]
algorithm = "sshsig"
key_id = "kennel-maint-2026"
signature = "-----BEGIN SSH SIGNATURE-----\nU1NIU0lHAAAAAQAAADMAAAALc3NoLWVkMjU1MTkAAAAgKWVn2EUqe+ju7quWQv7aY/ihwM\npBFaipL7vU2UVH85IAAAAbcG9saWN5LnYxQHByb2plY3RrZW5uZWwub3JnAAAAAAAAAAZz\naGE1MTIAAABTAAAAC3NzaC1lZDI1NTE5AAAAQO4+Qm6EihZIVrHHJHc5lYaTh22kJbwk+8\nYVUolVJ+bBxwFC1WisLKlFSZ4fXR916/2Lo8NXG0LCGg+pvSOcygQ=\n-----END SSH SIGNATURE-----\n"
\end{kennelcode}

\hypertarget{spawn}{%
\section{\texorpdfstring{\texttt{{[}spawn{]}}}{{[}spawn{]}}}\label{spawn}}

\texttt{{[}spawn{]}}: the delegated-instantiation grant.

\begin{kennelcode}
[spawn]
# optional
max_instances = 0
# decoration
reason   = "..."
\end{kennelcode}

Contains: \texttt{{[}{[}spawn.allow{]}{]}}.

\begin{description}
\tightlist
\item[\texttt{max\_instances}]
optional. Concurrent-instance ceiling across this grant's spawns, the
fork-bomb bound.
\item[\texttt{reason}]
decoration. Why this delegation is extended (required; the spawn waiver
is loud, validated at compile).
\end{description}

\textbf{\texttt{{[}{[}spawn.allow{]}{]}}} entries, optional:

\begin{description}
\tightlist
\item[\texttt{mutable}]
optional. Optional per-requester narrowing: the subset of the template's
\texttt{{[}{[}mutable{]}{]}} manifest fields
\item[\texttt{template}]
optional. The trust-store template name (\texttt{net-fetch}).
\end{description}

\hypertarget{ssh}{%
\section{\texorpdfstring{\texttt{{[}ssh{]}}}{{[}ssh{]}}}\label{ssh}}

\texttt{{[}ssh{]}}: per-kennel SSH egress (source-only).

\begin{kennelcode}
[ssh]
# optional
allow_headless = true
\end{kennelcode}

Contains: \texttt{{[}{[}ssh.destinations{]}{]}}.

\begin{description}
\tightlist
\item[\texttt{allow\_headless}]
optional. Whether a granted key may be driven by a non-interactive (CI)
kennel with no
\end{description}

\textbf{\texttt{{[}{[}ssh.destinations{]}{]}}} entries (or
\texttt{\{\ add,\ remove\ \}} increment), optional:

\begin{description}
\tightlist
\item[\texttt{dest}]
optional. The SSH destination, in the form the host-side \texttt{ssh} is
invoked with
\item[\texttt{options}]
optional. Host-side \texttt{ssh} invocation options for this
destination, passed verbatim as argv
\item[\texttt{reason}]
decoration. Why this destination is granted (required).
\item[\texttt{threats}]
decoration. Threat tags.
\end{description}

\hypertarget{trust}{%
\section{\texorpdfstring{\texttt{{[}trust{]}}}{{[}trust{]}}}\label{trust}}

\texttt{{[}trust{]}}: the masked workspace manifest (T2.8).

\begin{kennelcode}
[trust]
# optional
manifest = true
on_change = "warn" | "freeze" | "kill"
\end{kennelcode}

\begin{description}
\tightlist
\item[\texttt{manifest}]
optional. Maintain a \texttt{.trust-manifest.json} at the root of every
writable/persistent
\item[\texttt{on\_change}]
optional. What \texttt{kenneld} does when a watched trigger is mutated
during the run: One of \texttt{warn}, \texttt{freeze}, \texttt{kill}.
\end{description}

\hypertarget{tty}{%
\section{\texorpdfstring{\texttt{{[}tty{]}}}{{[}tty{]}}}\label{tty}}

\texttt{{[}tty{]}}: terminal hardening for an interactive (PTY)
workload.

\begin{kennelcode}
[tty]
# optional
filter_terminal_escapes = true
\end{kennelcode}

\begin{description}
\tightlist
\item[\texttt{filter\_terminal\_escapes}]
optional. Filter the dangerous escape sequences a workload could write
toward the
\end{description}

\hypertarget{unix}{%
\section{\texorpdfstring{\texttt{{[}unix{]}}}{{[}unix{]}}}\label{unix}}

\texttt{{[}unix{]}}: \texttt{AF\_UNIX} policy.

\begin{kennelcode}
[unix]
# optional
abstract = "deny" | "allow"
\end{kennelcode}

Contains: \texttt{{[}{[}unix.allow{]}{]}}.

\begin{description}
\tightlist
\item[\texttt{abstract}]
optional. Abstract-namespace socket disposition (\texttt{"deny"} /
\texttt{"allow"}). One of \texttt{deny}, \texttt{allow}.
\end{description}

\textbf{\texttt{{[}{[}unix.allow{]}{]}}} entries (or
\texttt{\{\ add,\ remove\ \}} increment), optional:

\begin{description}
\tightlist
\item[\texttt{env}]
optional. An environment variable to set to the shim path
(e.g.~\texttt{SSH\_AUTH\_SOCK}).
\item[\texttt{name}]
optional. A logical name (e.g.~\texttt{"ssh-agent"}) for a per-kennel
service instance.
\item[\texttt{real}]
optional. The real host socket path.
\item[\texttt{reason}]
decoration. Why this socket is granted (required).
\item[\texttt{shim}]
optional. The shim path the socket is bound at inside the kennel.
\item[\texttt{threats}]
decoration. Threat tags.
\end{description}

Example, from \texttt{toml/templates/dbus-broker}:

\begin{kennelcode}
[unix]
\end{kennelcode}

\hypertarget{unsafe}{%
\section{\texorpdfstring{\texttt{{[}unsafe{]}}}{{[}unsafe{]}}}\label{unsafe}}

\texttt{{[}unsafe{]}}: the advisory footgun umbrella.

Contains: \texttt{{[}unsafe.ptrace{]}}, \texttt{{[}unsafe.signal{]}}.

\hypertarget{unsafe.ptrace}{%
\subsection{\texorpdfstring{\texttt{{[}unsafe.ptrace{]}}}{{[}unsafe.ptrace{]}}}\label{unsafe.ptrace}}

\texttt{{[}unsafe.ptrace{]}}: ptrace across the kennel boundary (scoping
from PID-ns + seccomp).

\begin{kennelcode}
[unsafe.ptrace]
# optional
allow_from = ["..."]
allow_targets = ["..."]
\end{kennelcode}

\begin{description}
\tightlist
\item[\texttt{allow\_from}]
optional. Permitted sources.
\item[\texttt{allow\_targets}]
optional. Permitted targets (\texttt{"self"}, \ldots).
\end{description}

\hypertarget{workload}{%
\section{\texorpdfstring{\texttt{{[}workload{]}}}{{[}workload{]}}}\label{workload}}

\texttt{{[}workload{]}}: the command the kennel runs, optionally pinned.

\begin{kennelcode}
[workload]
# optional
argv     = ["..."]
cwd      = "..."
pinned   = true
sha256   = ["..."]
\end{kennelcode}

\begin{description}
\tightlist
\item[\texttt{argv}]
optional. The command + args (\texttt{argv{[}0{]}} is the program).
Absent ⇒ supplied at \texttt{kennel\ run}.
\item[\texttt{cwd}]
optional. Working directory inside the view (may carry a
\texttt{\textasciitilde{}}/\texttt{\textless{}home\textgreater{}}
placeholder).
\item[\texttt{pinned}]
optional. Refuse a CLI \texttt{-\/-} override of \texttt{argv} unless
\texttt{-\/-force} (pin exactly what runs).
\item[\texttt{sha256}]
optional. Accepted lowercase-hex SHA-256 digests of the workload binary;
the spawn verifies
\end{description}

Example, from \texttt{toml/policies/interactive}:

\begin{kennelcode}
[workload]
argv = ["/bin/bash"]
\end{kennelcode}

\hypertarget{common-forms}{%
\section{Common forms}\label{common-forms}}

\textbf{Threat tags}: carried inline on any grant that can widen
surface, as
\texttt{threats\ =\ \{\ exposed\ =\ {[}...{]},\ mitigated\ =\ {[}...{]}\ \}}.

\textbf{The increment}: any list field also accepts
\texttt{{[}{[}\textless{}list\textgreater{}.add{]}{]}} /
\texttt{{[}{[}\textless{}list\textgreater{}.remove{]}{]}}, folded over
the inherited list at compile.
